% Lame de zinc dans une solution de sulfate de cuivre : observations et interprétation (L'essentiel, chapitre 3)
\begin{tikzpicture}[x=1cm,y=1cm]
% Avant
\pic at (0,0) {becher={largeur=1.8,hauteur=2.2,niveau=1.6,liquide=cyan!40!blue!35}};
\filldraw[fill=otMuted!45,draw=otInk!70,line width=0.5pt] (-0.12,0.3) rectangle (0.12,2.7);
\node[pcLegendeG,anchor=north] at (0,-0.15) {au début};
\node[pcLegende,anchor=north,align=center] at (0,-0.55) {solution bleue\\(ions $\ce{Cu^{2+}}$)};
\draw[otGreenDark,line width=1pt,-{Stealth[length=2.5mm]}] (1.3,1.2) -- (2.3,1.2) node[midway,above,pcLegende,text=otGreenDark] {temps};
% Après
\pic at (3.6,0) {becher={largeur=1.8,hauteur=2.2,niveau=1.6,liquide=cyan!25!blue!12}};
\filldraw[fill=otMuted!45,draw=otInk!70,line width=0.5pt] (3.48,0.3) rectangle (3.72,2.7);
\fill[orange!70!red!80!black] (3.42,0.35) rectangle (3.48,1.55) (3.72,0.35) rectangle (3.78,1.55);
\foreach \y in {0.5,0.8,1.1,1.4}{\fill[orange!70!red!80!black] (3.40,\y) circle (0.05) (3.80,\y+0.1) circle (0.05);}
\node[pcLegendeG,anchor=north] at (3.6,-0.15) {après quelques minutes};
\node[pcLegende,anchor=north,align=center] at (3.6,-0.55) {dépôt rouge de cuivre ;\\solution décolorée};
% Interprétation microscopique
\begin{scope}[shift={(7.4,1.2)}]
  \fill[otMuted!40] (0,0) circle (0.42); \node[font=\otSemi\fontsize{8}{9}\selectfont] at (0,0) {Zn};
  \draw[otInk!70,line width=0.7pt,-{Stealth[length=2mm]}] (0.5,0) -- (1.3,0);
  \fill[otMuted!20] (1.75,0) circle (0.36); \node[font=\otSemi\fontsize{7.5}{9}\selectfont] at (1.75,0) {Zn$^{2+}$};
  \fill[cyan!35] (4.3,0) circle (0.36); \node[font=\otSemi\fontsize{7.5}{9}\selectfont] at (4.3,0) {Cu$^{2+}$};
  \draw[otInk!70,line width=0.7pt,-{Stealth[length=2mm]}] (4.75,0) -- (5.45,0);
  \fill[orange!70!red!80!black] (5.85,0) circle (0.38); \node[font=\otSemi\fontsize{8}{9}\selectfont,text=white] at (5.85,0) {Cu};
  \draw[otGreenDark,line width=1.1pt,-{Stealth[length=2.4mm]}] (0.3,0.4) to[out=50,in=130] node[midway,above,pcLegende,text=otGreenDark] {2 électrons transférés} (4.1,0.4);
  \node[pcLegende,anchor=north,align=center] at (0.85,-0.55) {oxydation :\\$\ce{Zn -> Zn^{2+} + 2e-}$};
  \node[pcLegende,anchor=north,align=center] at (5.0,-0.55) {réduction :\\$\ce{Cu^{2+} + 2e- -> Cu}$};
\end{scope}
\end{tikzpicture}
\par\smallskip\legende{Le zinc cède des électrons aux ions cuivre(II) : il s'oxyde en ions zinc(II), tandis que les ions cuivre(II) sont réduits en cuivre métallique.}
